\begin{explanationbox}
	\textbf{\textcolor{CExplanation}{一句话直觉}}

	绝对值对应数轴上点到原点的距离，因此绝对值不等式直接反映距离与数的大小关系。

	\medskip
	\textbf{\textcolor{CExplanation}{核心拆解}}
	\begin{description}[style=nextline, leftmargin=2.8em, labelsep=0.5em, itemsep=0.45em]
		\item[\textbf{要点一（几何意义）：}] $|f(x)|<a$ 表示 $f(x)$ 对应的点到原点的距离小于 $a$，所以 $f(x)$ 限制在 $(-a,a)$ 内。
		\item[\textbf{要点二（逻辑关系）：}] $|f(x)|>a$ 表示距离大于 $a$，所以 $f(x)$ 必须在 $(-a,a)$ 之外，即 $f(x)<-a$ 或 $f(x)>a$。
	\end{description}

	\medskip
	\textbf{\textcolor{CExplanation}{几何本质（数轴距离）}}

	$|f(x)|$ 是数轴上点 $f(x)$ 到原点的距离。 不等式 $|f(x)|<a$ 的几何意义是点 $f(x)$ 到原点的距离小于 $a$，故该点位于 $(-a,a)$ 内。 同理 $|f(x)|>a$ 表示点位于区间之外。

	\medskip
	\textbf{\textcolor{CExplanation}{代数意义（分段转化）}}

	通过绝对值的定义，将 $|f(x)|$ 按 $f(x)\ge 0$ 和 $f(x)<0$ 分段，得到两个等价的不等式，再合并成结论中的形式。

	\medskip
	\textbf{\textcolor{CExplanation}{考点价值（解不等式核心）}}
	\begin{itemize}[leftmargin=*, itemsep=0.5em, topsep=0.4em]
		\item \textbf{考法一（直接求解）：} 给定具体绝对值不等式，直接套用等价形式转化为普通不等式（组）求解。
		\item \textbf{考法二（参数问题）：} 已知不等式的解集，反求参数 $a$ 的值或范围。
	\end{itemize}

	\medskip
	\textbf{\textcolor{CExplanation}{顿悟点}}

	关键是把“绝对值符号”转化为“区间范围”，而且注意“小于取中间，大于取两边”的口诀。

	\medskip
	\textbf{\textcolor{CExplanation}{使用场景}}
	\begin{itemize}[leftmargin=*, itemsep=0.5em, topsep=0.4em]
		\item \textbf{场景一（单绝对值不等式）：} 含一个绝对值表达式的不等式。
		\item \textbf{场景二（含参绝对值不等式）：} 系数中含参数，需分类讨论。
	\end{itemize}
\end{explanationbox}
